\item \textbf{Diets.}  Developers of diets like to advertise that their diets produce significant weight loss results.  A new diet has been developed and several studies were done to investigate the effectiveness of the diet. In each study the average number of pounds lost was of interest, and researchers would like to show that dieters lost a significant amount of weight.
\begin{enumerate}
\item  Identify $\mu$ in this scenario. \textbf{(Choose One)}
\begin{itemize}
	\item Average number of pounds lost by all people in the sample.
	\item True average weight of all people in the population.
	\item \red{True average number of pounds lost by all people in the population.}
	\item Average weight of all people in the sample.
\end{itemize}
\item   Identify the appropriate null and the alternative hypotheses to test the researchers' claim.
$$\red{H_0: \mu = 0 \qquad \mbox{versus}  \qquad  H_a: \mu > 0}$$
\item  Suppose that an average weight loss of 1 pound (${\overline x} = 1$) and standard deviation of $s = 7$ pounds were found in each of the following scenarios. Further, assume weight loss is normally distributed. For each case, compute the test statistic, find the approximate $p$-value and make a decision about the null hypothesis using an $\alpha = 0.05$ level of significance. Round $t$-statistics to two decimals.
\begin{enumerate}
\item  \textbf{\underline{Scenario 1:}} a sample of size $n = 10$\\
\begin{enumerate}
\item  $$\red{t = \frac{\bar{x} - \mu_0}{s/\sqrt{n}} = \frac{1 - 0}{7/\sqrt{10}} = \mathbf{0.45} \qquad \qquad df = \mathbf{9}
}$$

\item The $p$-value corresponding to the correct test statistic is 0.331. This $p$-value suggests \underline{\red{Little or no evidence}} against the null hypothesis. 

%\begin{itemize}
%\item \red{Little or no evidence}
%\item Weak evidence
%\item Moderately strong evidence
%\item Strong evidence
%\item Extremely strong evidence
%\end{itemize}
%\item[]
\item \underline{Decision:}
\begin{itemize}
\item Reject the null hypothesis because the $p$-value $\leq$ $\alpha$.
\item Reject the null hypothesis because the $p$-value $>$ $\alpha$.
\item Fail to reject the null hypothesis because the $p$-value $\leq$ $\alpha$.
\item \red{Fail to reject the null hypothesis because the $p$-value $>$ $\alpha$.}
\end{itemize}
\end{enumerate}
\item[]
\item \textbf{\underline{Scenario 2:}} a sample of size $n = 120$
\begin{enumerate}
\item  $$\red{t = \frac{\bar{x} - \mu_0}{s/\sqrt{n}} = \frac{1 - 0}{7/\sqrt{120}} = \mathbf{1.56} \qquad \qquad df = \mathbf{119}}$$

\item  The $p$-value corresponding to the correct test statistic is 0.0601. This $p$-value suggests \underline{\red{Moderately strong evidence}} against the null hypothesis.

%\begin{itemize}
%\item Little or no evidence
%\item Weak evidence
%\item \red{Moderately strong evidence}
%\item Strong evidence
%\item Extremely strong evidence
%\end{itemize}
%\item[]
\item  \underline{Decision:}
\begin{itemize}
\item Reject the null hypothesis because the $p$-value $\leq$ $\alpha$.
\item Reject the null hypothesis because the $p$-value $>$ $\alpha$.
\item Fail to reject the null hypothesis because the $p$-value $\leq$ $\alpha$.
\item \red{Fail to reject the null hypothesis because the $p$-value $>$ $\alpha$.}
\end{itemize}
\end{enumerate}
\item[]
\item  \textbf{\underline{Scenario 3:}} a sample of size $n = 1050$\\ 
\begin{enumerate}
\item $$\red{t = \frac{\bar{x} - \mu_0}{s/\sqrt{n}} = \frac{1 - 0}{7/\sqrt{1050}} = \mathbf{4.63} \qquad \qquad df = \mathbf{1049}}$$

\item   The $p$-value corresponding to the correct test statistic is 0.000002. This $p$-value suggests \underline{\red{Extremely strong evidence}} against the null hypothesis.
%\begin{itemize}
%\item Little or no evidence
%\item Weak evidence
%\item Moderately strong evidence
%\item Strong evidence
%\item \red{Extremely strong evidence}
%\end{{itemize}
%\item[]
\item  \underline{Decision:}
\begin{itemize}
\item \red{Reject the null hypothesis because the $p$-value $\leq$ $\alpha$.}
\item Reject the null hypothesis because the $p$-value $>$ $\alpha$.
\item Fail to reject the null hypothesis because the $p$-value $\leq$ $\alpha$.
\item Fail to reject the null hypothesis because the $p$-value $>$ $\alpha$.
\end{itemize}
%\end{enumerate}
%\end{enumerate}
\item[]
\item   A weight loss of one pound is found to be statistically significant. The researchers share their results with the public claiming that their diet produces meaningful, significant weight loss. Do you agree? Explain your answer.\\
\emph{\red{From a practical standpoint, losing one pound is not a significant amount of weight}}.

\end{enumerate}
%\item Based on the results of your three tests, what can you say about the effect of sample size %on producing statistically significant results?\\
%\emph{Looking at the results from above, we can see that increasing the size of the sample taken %can make even small differences statistically significant.}
\end{enumerate}
\end{enumerate}